% Full editable typesetting of the retained Markdown source.
% Source: NILRADICAL_COTANGENT_DIVIDED_RESIDUE_DERIVATION.md
% Source SHA256: 38f5345eb82d6cca9975c76be63f40e0813b58239db2f654098d4da5d9fe37a9
% This conversion is not a new proof or PDF build.
% Options for packages loaded elsewhere
\PassOptionsToPackage{unicode}{hyperref}
\PassOptionsToPackage{hyphens}{url}
\documentclass[
  10pt,
]{article}
\usepackage{xcolor}
\usepackage[margin=22mm]{geometry}
\usepackage{amsmath,amssymb}
\setcounter{secnumdepth}{-\maxdimen} % remove section numbering
\usepackage{iftex}
\ifPDFTeX
  \usepackage[T1]{fontenc}
  \usepackage[utf8]{inputenc}
  \usepackage{textcomp} % provide euro and other symbols
\else % if luatex or xetex
  \usepackage{unicode-math} % this also loads fontspec
  \defaultfontfeatures{Scale=MatchLowercase}
  \defaultfontfeatures[\rmfamily]{Ligatures=TeX,Scale=1}
\fi
\usepackage{lmodern}
\ifPDFTeX\else
  % xetex/luatex font selection
\fi
% Use upquote if available, for straight quotes in verbatim environments
\IfFileExists{upquote.sty}{\usepackage{upquote}}{}
\IfFileExists{microtype.sty}{% use microtype if available
  \usepackage[]{microtype}
  \UseMicrotypeSet[protrusion]{basicmath} % disable protrusion for tt fonts
}{}
\makeatletter
\@ifundefined{KOMAClassName}{% if non-KOMA class
  \IfFileExists{parskip.sty}{%
    \usepackage{parskip}
  }{% else
    \setlength{\parindent}{0pt}
    \setlength{\parskip}{6pt plus 2pt minus 1pt}}
}{% if KOMA class
  \KOMAoptions{parskip=half}}
\makeatother
\setlength{\emergencystretch}{3em} % prevent overfull lines
\providecommand{\tightlist}{%
  \setlength{\itemsep}{0pt}\setlength{\parskip}{0pt}}
\usepackage{mathrsfs}
\usepackage{mathtools}
\usepackage{xurl}
\emergencystretch=3em
\setcounter{tocdepth}{1}
\newcommand{\Beta}{\mathrm{B}}
\DeclareUnicodeCharacter{00B2}{\ensuremath{^2}}
\DeclareUnicodeCharacter{00B1}{\ensuremath{\pm}}
\DeclareUnicodeCharacter{00D7}{\ensuremath{\times}}
\DeclareUnicodeCharacter{2192}{\ensuremath{\to}}
\DeclareUnicodeCharacter{21A6}{\ensuremath{\mapsto}}
\DeclareUnicodeCharacter{2208}{\ensuremath{\in}}
\DeclareUnicodeCharacter{2209}{\ensuremath{\notin}}
\DeclareUnicodeCharacter{2212}{\ensuremath{-}}
\DeclareUnicodeCharacter{2227}{\ensuremath{\wedge}}
\DeclareUnicodeCharacter{2228}{\ensuremath{\vee}}
\DeclareUnicodeCharacter{2260}{\ensuremath{\ne}}
\DeclareUnicodeCharacter{2264}{\ensuremath{\le}}
\DeclareUnicodeCharacter{2265}{\ensuremath{\ge}}
\DeclareUnicodeCharacter{2297}{\ensuremath{\otimes}}
\usepackage{bookmark}
\IfFileExists{xurl.sty}{\usepackage{xurl}}{} % add URL line breaks if available
\urlstyle{same}
\hypersetup{
  hidelinks,
  pdfcreator={LaTeX via pandoc}}

\author{}
\date{}

\begin{document}

\section{The nilradical--cotangent bridge and its divided residue at
3}\label{the-nilradicalcotangent-bridge-and-its-divided-residue-at-3}

Independent bounded derivation, 23 September 2026. This note proves the
exact module maps between the original nilradicals and cotangent
cohomology groups, keeps the integral quotient before and after
tensoring, and computes the divided residue on both cokernels. Division
by 3 occurs only in free lattices where the displayed coefficients prove
divisibility.

The retained source presentations and earlier computations are
\url{INTEGRAL_EIGHT_STATE_COLLISION_DERIVATION.md}, IEC41--IEC47 and
Appendix A, \url{IEC_P3_INDEPENDENT_REVIEW.md}, R3.1--R3.5, and
\url{INTEGRAL_RESIDUE_FROBENIUS_REVIEW.md}, RF39--RF54. The current maps
are proved in full below. No new publication or historical-literature
claim is made.

\subsection{NC1. The original algebras and their
radicals}\label{nc1.-the-original-algebras-and-their-radicals}

Put \(O=\mathbb Z_3\), and retain \[
B=O[x,T]/(f,g),\quad
f=x^2(x+3),\quad g=T^2-3x(x+2),
\] \[
C=O[\epsilon,T]/(f_C,g_C),\quad
f_C=\epsilon^2,\quad g_C=T^2-6\epsilon.
\tag{NC1}
\] The quotient \(q:B\to C\) sends \(x\mapsto\epsilon,T\mapsto T\). The
bases \[
1,x,x^2,T,xT,x^2T,\qquad
1,\epsilon,T,\epsilon T
\tag{NC2}
\] are free \(O\)-bases by successive division by the two monic
equations. Substitution verifies \(q\), and its images contain the basis
of \(C\), so it is surjective.

The nilradical of \(C\) is \[
N_C=O\epsilon\oplus OT\oplus O\epsilon T.
\] It has fourth power zero, and \(C/N_C=O\) is reduced, proving the
assertion. The nilradical of \(B\) has basis \[
n_1=x(x+3),\qquad n_2=T(x+3),\qquad n_3=Tx(x+3).
\tag{NC3}
\] For completeness, evaluate \(B\) at \((x,T)=(0,0),(-3,3),(-3,-3)\).
For an element \(p(x)+Tq_0(x)\), where both polynomials have degree at
most two, simultaneous vanishing means \[
p(0)=p(-3)=q_0(-3)=0.
\] This follows by adding and subtracting the last two evaluations;
cancellation of the nonzero scalar 3 is valid in the free coefficient
ring, without making it a unit. Coefficient comparison then gives \[
p=a(x^2+3x),\qquad
q_0=b(x+3)+c(x^2+3x).
\] Thus the common kernel is exactly the span of (NC3). Multiplication
gives \[
n_1^2=n_1n_3=n_2n_3=n_3^2=0,\quad
n_2^2=18n_1,\quad n_1n_2=3n_3.
\] The span is therefore nilpotent; every nilpotent must vanish in the
reduced evaluation target \(O^3\). This proves it is precisely \(N_B\).

The complete actions needed below are \[
xn_1=xn_3=0,\quad xn_2=n_3,\qquad
Tn_1=n_3,\quad Tn_2=6n_1,\quad Tn_3=0,
\tag{NC4}
\] and \[
\epsilon^2=0,\quad \epsilon T=\epsilon T,\quad
\epsilon(\epsilon T)=0,\qquad
T\epsilon=\epsilon T,\quad T^2=6\epsilon,\quad T(\epsilon T)=0
\tag{NC5}
\] on the indicated bases. For example \(Tn_2=3x(x+2)(x+3)=6n_1\), using
\(xn_1=0\). These identities prove (NC4); (NC5) is the presentation of
\(C\).

The radical map is \[
q_N(n_1)=3\epsilon,\qquad
q_N(n_2)=3T+\epsilon T,\qquad
q_N(n_3)=3\epsilon T.
\tag{NC6}
\] Its kernel is zero by coefficient comparison. Its image is a
\(C\)-submodule of \(N_C\), since \(q\) is onto and \(N_B\) is a
\(B\)-ideal. Its exact cokernel is \[
Q_N:=N_C/q_N(N_B)
=(O/3O)a\oplus(O/9O)b,\qquad
a=[\epsilon],\quad b=[T],\quad[\epsilon T]=-3b.
\tag{NC7}
\] Indeed the three relations in (NC6) are \(3\epsilon=0\),
\(\epsilon T=-3T\), and \(3\epsilon T=0\), which reduce exactly to
\(3a=0,9b=0\).

This is also a specified \(C\)-module, not only an \(O\)-module: \[
\epsilon a=0,\qquad \epsilon b=-3b,\qquad
Ta=-3b,\qquad Tb=0.
\tag{NC8}
\] The last identity uses \(6a=0\); all the others follow from (NC5) and
(NC7).

\subsection{\texorpdfstring{NC2. Cotangent complexes, their kernels, and
the map induced by
\(q\)}{NC2. Cotangent complexes, their kernels, and the map induced by q}}\label{nc2.-cotangent-complexes-their-kernels-and-the-map-induced-by-q}

Both pairs of defining equations are regular sequences in their
polynomial rings. The first equation is a nonzero divisor; modulo it,
the second is monic in \(T\), so multiplying a nonzero polynomial by it
retains a nonzero leading coefficient. The regular-immersion cotangent
presentations, in degrees \(-1,0\), are consequently \[
L_{B/O}=
\left[B\bar f\oplus B\bar g
\xrightarrow{J_B}B\,dx\oplus B\,dT\right],
\qquad
J_B=
\begin{pmatrix}3x(x+2)&-6(x+1)\\0&2T\end{pmatrix},
\] \[
L_{C/O}=
\left[C\bar f_C\oplus C\bar g_C
\xrightarrow{J_C}C\,d\epsilon\oplus C\,dT\right],
\qquad
J_C=
\begin{pmatrix}2\epsilon&-6\\0&2T\end{pmatrix}.
\tag{NC9}
\] Write \(H_B=H^{-1}(L_{B/O})\) and \(H_C=H^{-1}(L_{C/O})\).

The free \(O\)-bases are \[
u_1=(n_1,0),\qquad u_2=(n_3,0),\qquad u_3=(n_2,n_3)
\tag{NC10}
\] for \(H_B\), and \[
v_1=(\epsilon,0),\qquad v_2=(\epsilon T,0),\qquad
v_3=(3T,\epsilon T)
\tag{NC11}
\] for \(H_C\).

Here is a direct kernel proof. The ring \(B\) is local: its reduction
modulo 3 is \(\mathbb F_3[x,T]/(x^3,T^2)\), and a maximal ideal in the
finite integral \(O\)-algebra \(B\) lies over the maximal ideal of
\(O\). Thus \(x+1,x+2\) are units. In \(O[x]/f\), multiplying
\(a+bx+cx^2\) by \(x\) gives \(ax+(b-3c)x^2\), so its kernel is
\(On_1\). Expanding \(p+Tq_0\) now proves \[
\operatorname{ann}_B(T)=On_3,\qquad
\operatorname{ann}_B(T^2)=On_1\oplus On_3.
\] The lower row of \(J_B(a,b)=0\) therefore gives \(b=cn_3\). Since
\((x+1)n_3=n_3\) and \(T^2n_2=6n_3\), the upper row gives
\(a-cn_2\in On_1\oplus On_3\). This proves (NC10) and its independence.

For \(C\), the equation \(Tb=0\) in its free basis gives
\(b=c\epsilon T\). The remaining equation \(\epsilon a=3b\) forces the
scalar coefficient of \(a\) to be zero and its \(T\)-coefficient to
equal \(3c\). The two other coefficients are unrestricted. This proves
(NC11). These arguments cancel nonzero scalars only in free modules.

The module actions on (NC10) are \[
xu_1=xu_2=0,\quad xu_3=u_2,\qquad
Tu_1=u_2,\quad Tu_2=0,\quad Tu_3=6u_1.
\tag{NC12}
\] The actions on (NC11) are \[
\epsilon v_1=\epsilon v_2=0,\quad\epsilon v_3=3v_2,\qquad
Tv_1=v_2,\quad Tv_2=0,\quad Tv_3=18v_1.
\tag{NC13}
\] They follow by multiplying both coordinates of each displayed cycle.

The polynomial lift of \(q\) gives the relation map \[
\bar f\longmapsto(\epsilon+3)\bar f_C,\qquad
\bar g\longmapsto-3\bar f_C+\bar g_C,
\tag{NC14}
\] because \(f(\epsilon)=(\epsilon+3)f_C\) and
\(g(\epsilon,T)=g_C-3f_C\). Its differential map is
\(dx\mapsto d\epsilon,dT\mapsto dT\). Differentiating these two
polynomial identities proves the chain identity. The induced map is \[
q_Hu_1=9v_1,\qquad q_Hu_2=9v_2,\qquad
q_Hu_3=3v_3-3v_2.
\tag{NC15}
\] For the last coefficient, use \(q(n_2)=3T+\epsilon T\) and
\(q(n_3)=3\epsilon T\) in (NC14). The first component becomes
\(9T-3\epsilon T\), and the second \(3\epsilon T\), giving exactly
(NC15).

With \(w=v_3-v_2\), the map has diagonal entries \(9,9,3\). It is
injective and its cokernel is \[
Q_H:=H_C/q_H(H_B)
=(O/9O)c_1\oplus(O/9O)c_2\oplus(O/3O)w,
\quad c_1=[v_1],\quad c_2=[v_2].
\tag{NC16}
\] Its \(C\)-action is \[
\epsilon c_1=\epsilon c_2=0,\quad\epsilon w=3c_2,\qquad
Tc_1=c_2,\quad Tc_2=Tw=0.
\tag{NC17}
\] For \(Tw\), the free-lattice value is \(18v_1\), which is zero in
(NC16).

\subsection{NC3. The two radical--cotangent maps and their exact
compatibility}\label{nc3.-the-two-radicalcotangent-maps-and-their-exact-compatibility}

Define \[
\kappa_B:N_B\longrightarrow H_B,\qquad
n_1\longmapsto u_1,\quad n_2\longmapsto u_3,\quad n_3\longmapsto u_2,
\] \[
\kappa_C:N_C\longrightarrow H_C,\qquad
\epsilon\longmapsto3v_1,\quad
T\longmapsto v_3-2v_2,\quad
\epsilon T\longmapsto3v_2.
\tag{NC18}
\] The first map is a \(B\)-linear isomorphism: its basis correspondence
is bijective, and every action by \(x,T\) in (NC4) agrees with the
corresponding action in (NC12).

The second map is \(C\)-linear. For the nonzero generator actions, \[
\epsilon\kappa_C(T)=3v_2=\kappa_C(\epsilon T),\qquad
T\kappa_C(\epsilon)=3v_2=\kappa_C(T\epsilon),
\] \[
T\kappa_C(T)=18v_1=\kappa_C(6\epsilon).
\] All the zero generator actions in (NC5) give zero by (NC13). These
checks prove compatibility with both algebra generators, hence with
every scalar. It is injective: a relation \(3av_1+b(v_3-2v_2)+3cv_2=0\)
first gives \(b=0\), then \(a=c=0\), using freeness.

The square commutes exactly: \[
q_H\kappa_B=\kappa_Cq_N.
\tag{NC19}
\] On \(n_1,n_3\), both sides give \(9v_1,9v_2\). On \(n_2\), \[
\kappa_C(3T+\epsilon T)=3(v_3-2v_2)+3v_2=3v_3-3v_2.
\] This is (NC15).

The cokernel of \(\kappa_C\) is free of rank two over \(\mathbb F_3\),
but its \(T\)-action must be retained. Explicitly \[
R_C:=H_C/\kappa_C(N_C)\simeq C/(3,\epsilon)
\simeq\mathbb F_3[T]/T^2,
\tag{NC20}
\] as \(C\)-modules. The isomorphism sends \[
[v_1]\longmapsto1,\qquad[v_2]\longmapsto T,\qquad
[v_3]\longmapsto2T.
\] The relations \(3v_1=3v_2=0\) and \(v_3=2v_2\) prove its bijectivity
as an \(O\)-module. Equation (NC13) proves its \(C\)-linearity:
\(T[v_1]=[v_2]\), \(T[v_2]=0\), and \(\epsilon\) acts as zero.

Consequently there is an exact sequence of \(C\)-modules \[
0\longrightarrow Q_N\xrightarrow{\bar\kappa}Q_H
\xrightarrow{\psi}R_C\longrightarrow0.
\tag{NC21}
\] The maps are \[
\bar\kappa(a)=3c_1,\qquad
\bar\kappa(b)=w-c_2,
\] \[
\psi(c_1)=1,\qquad\psi(c_2)=T,\qquad\psi(w)=T.
\tag{NC22}
\] For injectivity, \(3ac_1+b(w-c_2)=0\) in (NC16) forces \(b=0\) modulo
9 from its \(c_2\)-component, and \(a=0\) modulo 3 from its
\(c_1\)-component. The map \(\psi\) is onto by its values on
\(c_1,c_2\), and its kernel has exactly the relations generated by
\(\bar\kappa(a),\bar\kappa(b)\); equivalently it is the quotient by the
image of \(\kappa_C\). This proves exactness. The \(C\)-linearity also
follows directly from (NC8), (NC17), and (NC22).

As an \(O\)-module sequence, (NC21) is not split. A lift of
\(1\in C/(3,\epsilon)\) must have \(c_1\)-coefficient congruent to 1
modulo 3. Its multiple by 3 therefore has a nonzero \(c_1\)-component
modulo 9. No lift is killed by 3, so an \(O\)-linear section cannot
exist. In particular a \(C\)-linear section cannot exist.

\subsection{NC4. Tensoring the source cotangent complex
first}\label{nc4.-tensoring-the-source-cotangent-complex-first}

Since the terms of \(L_{B/O}\) are free \(B\)-modules, its derived
tensor with \(C\) is represented by ordinary tensoring its displayed
complex. Write \[
H_{B,C}=H^{-1}(L_{B/O}\otimes_B^{\mathbf L}C).
\] Its differential matrix is \[
\begin{pmatrix}6\epsilon&-6(\epsilon+1)\\0&2T\end{pmatrix}.
\] Its free \(O\)-basis is \[
w_1=(\epsilon,0),\qquad
w_2=(\epsilon T,0),\qquad
w_3=(T,\epsilon T).
\tag{NC23}
\] The second kernel equation gives \(b=c\epsilon T\). The first, after
cancelling the nonzero scalar 6 in the free module \(C\), gives
\(\epsilon a=c\epsilon T\). Its coefficient comparison gives precisely
this basis.

There is a \(C\)-linear isomorphism \[
j_C:N_C\xrightarrow{\sim}H_{B,C},\qquad
\epsilon\mapsto w_1,\quad T\mapsto w_3,\quad\epsilon T\mapsto w_2.
\tag{NC24}
\] Indeed \(\epsilon w_3=w_2\), \(Tw_1=w_2\), \(Tw_3=6w_1\), and all the
remaining generator actions agree with (NC5).

The natural base-change map and the map induced by (NC14) are \[
H_B\longrightarrow H_{B,C}:\quad
u_1\mapsto3w_1,\quad u_2\mapsto3w_2,\quad u_3\mapsto3w_3+w_2,
\] \[
H_{B,C}\longrightarrow H_C:\quad
w_1\mapsto3v_1,\quad w_2\mapsto3v_2,\quad w_3\mapsto v_3-2v_2.
\tag{NC25}
\] Substitution proves all six values. Under \(\kappa_B,j_C\), the first
map is exactly \(q_N\); the second composed with \(j_C\) is exactly
\(\kappa_C\). Thus the two maps are injective, have respective cokernels
\(Q_N\) and \(R_C\), and their composite has cokernel \(Q_H\). This is
the actual source of the exact sequence (NC21).

In particular the cokernel after tensoring the source complex first is
\(R_C\), not \(Q_H\). It is \(\mathbb F_3^2\) as an \(O\)-module, with
the nontrivial \(T\)-action in (NC20). The original free-lattice map has
cokernel \((O/9)^2\oplus O/3\), as in (NC16).

There is an explicit conormal interpretation of \(R_C\). The ideal
\(J=\ker q=x^2B\) has free basis \(x^2,x^2T\), and \[
J^2=9J,\qquad J/J^2=(O/9O)[x^2]\oplus(O/9O)[x^2T].
\tag{NC26}
\] Indeed \(x^4=9x^2\) gives both generators of \(J^2\) and all their
products. Lifting \(f_C\) into \(B\) gives \(x^2\), while lifting
\(g_C\) gives \(3x^2\). With the connecting-map sign determined by these
lifts, the conormal map from the cycles is \[
\partial(v_1)=-3[x^2],\quad
\partial(v_2)=-3[x^2T],\quad
\partial(v_3)=-6[x^2T].
\tag{NC27}
\] For the first two values use \(x^3=-3x^2\); the third is
\(3Tx^2+3Tx^3=-6Tx^2\). A vector \(Av_1+Bv_2+Cv_3\) is in its kernel
exactly when \(A\in3O\) and \(B+2C\in3O\). This is precisely the image
lattice of \(\kappa_C\). Its image is \(3(J/J^2)\). Therefore \[
R_C\xrightarrow{\sim}3(J/J^2),\qquad
1\mapsto-3[x^2],\quad T\mapsto-3[x^2T].
\tag{NC28}
\] The \(C\)-action is the one in (NC20): \(\epsilon\) kills these
classes since multiplication by \(x\) supplies a further factor \(-3\),
and \(T\) carries the first to the second and kills the second. Thus
(NC28) is an exact \(C\)-module identification of the final term.

\subsection{NC5. The original residues and the factor
five}\label{nc5.-the-original-residues-and-the-factor-five}

Retain the derivations with their original coefficients: \[
D_B(x)=0,\quad D_B(T)=-n_3/8;\qquad
D_C(\epsilon)=0,\quad D_C(T)=-3\epsilon T/8.
\tag{NC29}
\] They are integral derivations. Their polynomial lifts satisfy \[
\widetilde D_B(f)=0,\qquad
\widetilde D_B(g)=-\frac34(x+2)f-\frac14n_1g,
\] \[
\widetilde D_C(f_C)=0,\qquad
\widetilde D_C(g_C)=-\frac92f_C-\frac34\epsilon g_C.
\tag{NC30}
\] For \(B\), expand \(D_B(g)=-n_1T^2/4\) and substitute
\(T^2=g+3x(x+2)\), using \(n_1x(x+2)=f(x+2)\). For \(C\), expand
\(-3\epsilon T^2/4\) and substitute \(T^2=g_C+6\epsilon\). This proves
preservation of the respective relation ideals. The maps commute with
\(q\) by their values on \(x,T\).

Their restrictions to the nilradicals are particularly exact: \[
D_B|_{N_B}=-\frac38\,x\cdot(-),\qquad
D_C|_{N_C}=-\frac38\,\epsilon\cdot(-).
\tag{NC31}
\] For \(B\), both sides kill \(n_1,n_3\) and send \(n_2\) to
\(-3n_3/8\), using \((x+3)n_3=3n_3\). For \(C\), both kill
\(\epsilon,\epsilon T\) and send \(T\) to \(-3\epsilon T/8\). Hence
these are module-linear maps on the displayed ideals, in addition to
being restrictions of derivations.

The polynomial relation formulas (NC30) induce Lie maps on the cotangent
complexes: on coefficients they satisfy
\(\mathcal L_D(am)=D(a)m+a\mathcal L_D(m)\), on relation generators they
use (NC30), and on differentials they send \(da\) to \(dD(a)\).
Differentiating (NC30) proves the chain identity. On the free cycle
bases their maps are \[
\mathcal L_Bu_1=\mathcal L_Bu_2=0,\qquad
\mathcal L_Bu_3=-15u_2/8,
\] \[
\mathcal L_Cv_1=\mathcal L_Cv_2=0,\qquad
\mathcal L_Cv_3=-45v_2/8.
\tag{NC32}
\] For \(u_3\), differentiating its coefficient \(n_2\) gives
\(-3n_3/8\), and multiplying its second coefficient \(n_3\) by the
\(\bar f\)-coefficient of \(\mathcal L_B\bar g\) gives
\(-3(x+2)n_3/4=-3n_3/2\). Their sum is \(-15n_3/8\). For \(v_3\), the
coefficient \(3T\) contributes \(-9\epsilon T/8\), while
\(\epsilon T\mathcal L_C\bar g_C\) contributes \(-9\epsilon T/2\). Their
sum is \(-45\epsilon T/8\).

In module notation, (NC12)--(NC13) give \[
\mathcal L_B=-\frac{15}{8}\,x\cdot(-)\quad\text{on }H_B,\qquad
\mathcal L_C=-\frac{15}{8}\,\epsilon\cdot(-)\quad\text{on }H_C.
\tag{NC33}
\] Therefore the precise comparison is \[
\mathcal L_B\kappa_B=5\kappa_BD_B|_{N_B},\qquad
\mathcal L_C\kappa_C=5\kappa_CD_C|_{N_C}.
\tag{NC34}
\] It also follows directly from the basis formulas, and retains the
numerical factor 5.

The two residue maps on the radicals commute with \(q_N\), and the two
Lie maps commute with \(q_H\). For the latter, on the only nonzero
column, \[
\mathcal L_Cq_H(u_3)
=\mathcal L_C(3v_3-3v_2)=-135v_2/8
=q_H(-15u_2/8).
\tag{NC35}
\] The other columns vanish.

\subsection{NC6. Dividing the free maps and computing both
cokernels}\label{nc6.-dividing-the-free-maps-and-computing-both-cokernels}

The four maps just proved are divisible by 3 on their free lattices.
Define \[
A_B=\frac{D_B|_{N_B}}3=-\frac{x}{8}\cdot(-),\qquad
A_C=\frac{D_C|_{N_C}}3=-\frac{\epsilon}{8}\cdot(-),
\] \[
L_B=\frac{\mathcal L_B}3=-\frac{5x}{8}\cdot(-),\qquad
L_C=\frac{\mathcal L_C}3=-\frac{5\epsilon}{8}\cdot(-).
\tag{NC36}
\] Every right side is an integral module map. It is the unique map
whose triple is the corresponding original map, because each lattice is
\(O\)-free. All four maps have square zero, by (NC4), \(\epsilon^2=0\),
and (NC12).

The compatibility squares survive this division. Thus \(A_C\) induces an
operator \(A_N\) on \(Q_N\), and \(L_C\) induces an operator \(A_H\) on
\(Q_H\). From (NC8) and (NC17), \[
A_N(a)=0,\qquad A_N(b)=\frac38b,
\] \[
A_H(c_1)=A_H(c_2)=0,\qquad
A_H(w)=-\frac{15}{8}c_2.
\tag{NC37}
\] The relation \([\epsilon T]=-3b\) gives the positive sign in the
first nonzero formula.

Both induced operators are nonzero and square-zero. Their exact modules
are \[
\begin{aligned}
\ker A_N&=(O/3O)a\oplus3(O/9O)b,\\
\operatorname{im}A_N&=3(O/9O)b\simeq\mathbb F_3,\\
\operatorname{coker}A_N&=(O/3O)a\oplus(O/3O)b;
\end{aligned}
\tag{NC38}
\] \[
\begin{aligned}
\ker A_H&=(O/9O)c_1\oplus(O/9O)c_2,\\
\operatorname{im}A_H&=3(O/9O)c_2\simeq\mathbb F_3,\\
\operatorname{coker}A_H&=(O/9O)c_1
\oplus(O/3O)c_2\oplus(O/3O)w.
\end{aligned}
\tag{NC39}
\] For (NC38), multiplication by \(3/8\) on an order-nine cyclic module
has kernel and image its order-three subgroup. For (NC39), the
\(w\)-summand has order three and maps injectively onto the order-three
subgroup generated by \(3c_2\), since \(-15/8\) has valuation one. These
observations prove every kernel, image, and cokernel; the operator
squares also follow directly from these formulas.

The unscaled operators on both \(Q_N\) and \(Q_H\) are zero: they are
the triples of (NC37), whose images have order three. Their vanishing
does not discard the exact nonzero maps obtained by division on the free
source and target lattices before passing to the quotient.

The factor-five comparison survives on the cokernels: \[
A_H\bar\kappa=5\bar\kappa A_N,\qquad \psi A_H=0.
\tag{NC40}
\] For example, \[
A_H\bar\kappa(b)=-15c_2/8,\qquad
5\bar\kappa A_N(b)=\frac{15}{8}(w-c_2)=-15c_2/8,
\] since \(3w=0\). Equivalently, the exact sequence with its specified
operators is \[
0\longrightarrow(Q_N,\,5A_N)
\xrightarrow{\bar\kappa}(Q_H,\,A_H)
\xrightarrow{\psi}(R_C,\,0)
\longrightarrow0.
\tag{NC41}
\] The factor 5 is part of this equation. The last operator is zero
because \(L_C\) is multiplication by \(-5\epsilon/8\), while
\(\epsilon\) kills \(R_C\).

There is a valid divided Lie map on the tensor-first free kernel as
well. In (NC23), the base-changed Lie action has \[
\mathcal L_{B,C}(w_1)=\mathcal L_{B,C}(w_2)=0,\qquad
\mathcal L_{B,C}(w_3)=-15w_2/8.
\tag{NC42}
\] Indeed the coefficient \(T\) contributes \(-3\epsilon T/8\), and the
coefficient \(\epsilon T\) times the specialized \(\bar f\)-coefficient
from (NC30) contributes \(-3\epsilon T/2\). Its divided map is
multiplication by \(-5\epsilon/8\). Thus (NC25) also gives compatible
divided maps before taking the tensor-first cokernel, and its induced
map on \(R_C\) is zero. This proves the distinction directly at the
level of the actual complexes and their kernels.

Finally, \(D_B/3\) does not define an integral derivation on the whole
algebra \(B\). It would send \[
T\longmapsto-\frac1{24}n_3
=-\frac1{24}x^2T-\frac18xT,
\tag{NC43}
\] whose coefficient of the independent basis vector \(x^2T\) does not
belong to \(O\). Its restriction to \(N_B\), in contrast, is the
integral multiplication map \(-x/8\) in (NC36). This specifies exactly
where division is allowed and where it is not.

\subsection{NC7. The complete bridge}\label{nc7.-the-complete-bridge}

The original diagram is \[
\begin{matrix}
N_B&\xrightarrow{q_N}&N_C\\
\downarrow\kappa_B&&\downarrow\kappa_C\\
H_B&\xrightarrow{q_H}&H_C,
\end{matrix}
\] with \(\kappa_B\) an isomorphism and all three other displayed maps
injective. Its cokernel sequence is (NC21), whose final module is
\(C/(3,\epsilon)\), including its nonzero \(T\)-action. Factoring the
bottom map through the tensor-first kernel gives (NC25); this accounts
for the two different cotangent cokernels without replacing either one.

The original residues and their Lie actions are related by the exact
factor 5 in (NC34). Their free-lattice divisions give the nonzero
square-zero cokernel maps in (NC37), with complete kernels and images in
(NC38)--(NC39), and the exact operator sequence (NC41). The final module
also identifies with the conormal torsion subgroup \(3(J/J^2)\) by
(NC28). These are maps between the retained original objects with their
full integral scalar actions.

\end{document}
